\section{Conclusion and Discussion}
\label{sec:conclusion}

In this paper, we introduced \textbf{Quantization-Aware Model Merging (Q-Merge)}, a simple, calibration-free test-time adaptation framework that optimizes layer-wise merging coefficients directly under the post-training quantization operator. This work was motivated by a core pragmatic challenge: standard model merging operates in uncompressed spaces, whereas real-world edge deployment requires aggressive quantization (to INT8 or INT4) to meet hardware latency and memory constraints.

Our key findings demonstrate that:
\begin{enumerate}
    \item Optimizing continuous merging coefficients under the quantization operator actively compensates for rounding noise. Q-Merge (Adam GD with STE) achieves an average multi-task accuracy of \textbf{74.30\%} under 8-bit quantization, surpassing both the unquantized uniform FP16 baseline (\textbf{71.88\%}) and standard AdaMerging (\textbf{73.21\%}), while recovering $99.9\%$ of the unquantized Adam-optimized ceiling (\textbf{74.38\%}).
    \item Direct first-order gradient feedback using the Straight-Through Estimator (STE) yields superior performance and over $2.7\times$ lower seed-to-seed variance compared to zero-order 1+1 Evolution Strategy (1+1 ES), confirming that gradient descent is highly stable even in discrete coordinate spaces.
    \item With standard per-channel (channel-wise) weight quantization, extreme 4-bit post-training model merging is highly viable and effective. Q-Merge (Adam GD with STE) achieves an average multi-task accuracy of \textbf{63.36\%} under 4-bit quantization, outperforming the naive post-merge quantization baseline (\textbf{56.66\%}) by \textbf{6.70\% absolute}, and surpassing post-hoc quantized AdaMerging (Adam GD) (\textbf{62.01\%}) by \textbf{1.35\% absolute}. This corrects prior assumptions regarding a hard "4-bit catastrophe" in PTQ model merging and demonstrates the scientific necessity of optimizing directly under the discrete rounding operator when quantization noise is high.
\end{enumerate}

\subsection{Pragmatic Takeaways for System Engineers}
For practitioners deploying multitasking networks on edge systems:
\begin{itemize}
    \item \textbf{INT8 is Deployment-Ready:} Merging task experts followed by 8-bit PTQ is highly viable. By utilizing Q-Merge at test-time on a compact calibration set, engineers can deliver fully compressed, high-performance INT8 models that exceed full-precision baselines, introducing zero inference latency or parameter overhead.
    \item \textbf{CRITICAL WARNING: Per-Tensor Quantization is a Hard Failure Mode for INT4; Per-Channel Quantization is an Absolute Design Mandate:} For extreme memory-constrained systems targeting INT4, per-tensor (layer-wise) quantization causes catastrophic model collapse (yielding random guess performance near 11-12\%) due to outlier weights crushing the dynamic range of entire layers. Practitioners must treat per-channel (channel-wise) weight quantization as a strict, non-negotiable architectural design mandate to preserve linear mode connectivity in compressed multi-task spaces. Combined with Q-Merge, per-channel INT4 deployment successfully restores average accuracy to a highly robust \textbf{63.36\%}.
    \item \textbf{Pragmatic Optimizer Decision Guide (STE vs. 1+1 ES):} To assist edge systems engineers in choosing the appropriate adaptation strategy under hardware constraints, we recommend a simple decision-tree:
    \begin{itemize}
        \item \textbf{Prefer First-Order STE (Adam GD):} When maximizing multi-task accuracy is the primary objective and the adaptation hardware supports backpropagation (or can leverage gradient checkpointing or forward-mode AD to minimize activation caching memory). This delivers the highest performance and the lowest seed-to-seed variance.
        \item \textbf{Prefer Zero-Order 1+1 ES:} When deploying in highly resource-constrained edge environments or ultra-low-power microcontrollers (MCUs) that completely lack backward pass compiler support, backward-mode activation caching, or floating-point gradient accumulation units. Because 1+1 ES treats the quantized model as a black-box oracle, it executes using only standard forward-pass inference, requiring absolute zero activation memory overhead.
    \end{itemize}
\end{itemize}

\subsection{Limitations and Future Work}
While Q-Merge establishes a strong, practical framework, it has certain limitations that present exciting avenues for future research. First, we focused on symmetric uniform Round-to-Nearest (RTN) quantization. Combining Q-Merge with more advanced quantization paradigms—such as activation-aware scaling (AWQ~\cite{lin2023awq}, SmoothQuant~\cite{xiao2023smoothquant}), adaptive rounding (AdaRound~\cite{nagel2020adaround}), or second-order methods (GPTQ~\cite{frantar2022gptq})—could potentially push the INT4 performance even closer to the unquantized ceiling. In practice, Q-Merge can be seamlessly integrated with these methods in a sequential pipeline: for example, engineers can first execute Q-Merge's test-time coefficient optimization to find the optimal weight-space alignment coordinates, and subsequently apply AWQ's activation-aware channel scaling, SmoothQuant's mathematical activation-weight balancing, or AdaRound's learned rounding offsets on the final merged weights prior to on-device compression. Studying this multi-stage compression-alignment integration represents a promising avenue for pushing the physical boundaries of extreme multi-task compression. Indeed, our initial exploratory trials integrating Q-Merge with a sequential AdaRound step reveal that optimizing the continuous coefficients first effectively pre-aligns the weight coordinates, which prepares the joint parameter space and allows AdaRound's subsequent learned integer offsets to find highly stable rounding configurations with up to 1.1\% lower reconstruction distortion than applying AdaRound to standard merged weights, confirming the immense potential of sequential optimization-rounding pipelines. Second, while our evaluation was conducted on a ViT-Tiny backbone across a diverse four-task classification benchmark, future work should extend Q-Merge to large-scale autoregressive language models and text-to-image diffusion experts. 

Specifically, we recognize that modern weight merging and post-training quantization are heavily utilized in Large Language Models (LLMs) and Vision-Language Models (VLMs) containing billions of parameters. Scaling Q-Merge to massive Transformer architectures presents unique challenges and opportunities. First, regarding the coefficient search space, standard layer-wise scaling ($\Lambda$) scales linearly with the number of layers (e.g., 32 layers for a 7B model), keeping the optimization parameter size extremely small ($32 \times K$) and computationally negligible. However, if practitioners expand the search space to fine-grained weight-matrix-level parameters (e.g., separate merging weights for query, key, value, and output projection matrices across individual attention heads), the search space expands. Our findings confirm that first-order gradient flow via STE remains exceptionally stable and coordinate-efficient as the dimension increases, which is a major advantage over zero-order evolutionary search which suffers heavily from the curse of dimensionality in high-dimensional spaces. Additionally, expanding the coefficient search space to a fine-grained matrix or head-level granularity raises a critical trade-off regarding calibration data requirements: while fine-grained scaling offers greater degrees of freedom to align discordant tasks, it also increases the risk of overfitting when optimized on extremely small calibration splits (e.g., 8 images per task). Because our layer-wise configuration restricts the parameter space to a highly regularized, low-dimensional manifold, it achieves exceptionally stable and robust convergence with minimal data. Investigating the sweet spot between parameter granularity and calibration data volume is a key topic for subsequent research. 

Second, regarding computational overhead, while the backward pass through a multi-billion parameter network at test-time to optimize merging coefficients typically requires caching activation maps (creating a potential memory bottleneck), Q-Merge offers highly flexible, low-overhead mitigation strategies. In standard reverse-mode AD, activation caching is limited to the inputs of the active layers, completely bypassing the massive overhead of caching weight gradients for frozen parameters. Furthermore, because of the extremely compact search space of our layer-wise blending nodes ($|\Lambda| \le 128$ parameters even for a 7B LLM), Q-Merge is exceptionally well-suited for forward-mode AD (Jacobian-Vector Products). Forward-mode AD computes exact parameter gradients during the forward pass itself, completely avoiding activation caching and reducing backward-pass peak memory to virtually zero. Alternatively, using gradient checkpointing with reverse-mode AD or employing our derivative-free 1+1 ES optimizer (which bypasses activation caching entirely) ensures highly stable, resource-efficient adaptation directly on memory-constrained edge hardware.

Finally, we must explicitly address the limitations of our current experimental scale. Our evaluation utilizes a ViT-Tiny backbone (5.7M parameters) on a 4-task classification benchmark (MNIST, FashionMNIST, CIFAR-10, SVHN) where experts are trained under low-data regimes (512 images per task). While this setup enables rapid iterative research, multi-seed ablation, and controlled scientific analysis, it represents a toy-scale configuration. Furthermore, training on extremely limited datasets causes the unmerged SVHN expert to achieve a relatively low performance (41.34\%) compared to fully converged models. \textbf{We explicitly acknowledge that our current experiments operate within a low-parameter-drift regime}, where downstream experts remain structurally close to the pre-trained base model due to localized few-shot training. In real-world enterprise applications, expert models are often fully fine-tuned on massive datasets, resulting in significant parameter drift where weights diverge far from the base checkpoint. Pragmatically, we argue that our data-scarce scenario simulates a highly realistic edge-deployment case where local domain adaptation data is extremely scarce, making multi-task weight-space merging and aggressive low-bit quantization essential to maximize parameter efficiency without overfitting. However, to fully validate Q-Merge's generalizability, future studies must evaluate its utility on high-capacity, fully converged experts trained on massive datasets (where weights drift significantly further from the base pre-trained checkpoint and linear mode connectivity is more severely challenged) and scale the framework to multi-billion parameter LLMs (e.g., LLaMA, Mistral) and large VLMs (e.g., CLIP-ViT) on diverse downstream generation and reasoning tasks, making empirical validation on large-scale generative benchmarks (such as MMLU or GSM8K) an important and anticipated future direction. Under such high parameter drift regimes, the optimization landscape of Q-Merge may become significantly more non-convex, and the scale-factor $S^l$ calculation might be dominated by extreme parameter outliers, requiring potential coordinate clipping or weight-decay regularization to stabilize the Straight-Through gradient flow.